\section{Conclusion}
\label{sec:conclusion}

We formulated complex multi-hop retrieval as \emph{declarative evidence
execution}: a question is compiled into typed, dependent evidence requirements;
a deterministic importance-weighted allocator assigns legal orders, physical
implementations, and budget allocations to those requirements under a matched
budget; and a deterministic executor acquires evidence with hard budget
enforcement and full provenance. We introduced a typed evidence algebra whose
state-transition semantics (\code{apply\_observation}) make requirement
progress first-class and auditable, and we derived a chain law pinning
requirement importance to dependency depth ($\tau = 2d-1$) on well-defined
chains, so the estimator is construction-derived rather than learned.

The evaluation establishes a cost, not an accuracy, claim, and we are
deliberate about that boundary. Under the matched-budget protocol, on the
matched valid pairs the requirement-aware optimizer does not spend more on
average than the frozen static baseline:
realized retrieval calls fall on every dataset ($-0.19$ on HotpotQA, $-0.24$
on 2Wiki, $-0.09$ on MuSiQue; pooled $-0.16$ calls/question), though none of
these reductions is statistically supported (bootstrap $p \in [0.119,0.376]$;
Holm-adjusted $p$ up to $0.581$). EM differences are directionally positive
but equally non-significant ($\Delta$EM $+6.3/+5.9/+9.1$pt; McNemar
$p\geq 0.625$; pooled $+7.3$pt on 55 matched questions). On the $\ge 3$-slot
subdomain the cost saving is larger and statistically supported in the cleanest
comparison: pooled chain $3.45$ vs.\ cost-only $4.73$ calls/question
($\Delta=-1.27$, bootstrap $p<0.001$; vs.\ static $-1.09$, $p=0.036$) at pooled
EM that is equal across arms ($54.5\%/45.5\%/54.5\%$). No quality claim is made on
EM. The $1$--$2$-slot chains---which dominate most QA splits---leave the
optimizer little allocation freedom, so the cost saving is concentrated on the
$\ge 3$-slot subdomain.
We report the point-estimate quality difference as a point estimate, and we treat
runtime re-optimization over alternative physical plans as future work under a
branching execution model, because on the chain-only topology this system runs
today it is vacuous.

More broadly, this study suggests that complex retrieval can be treated as an
execution problem in which evidence goals remain stable while physical
acquisition strategies adapt to observed state and resource constraints. The
honest path forward is to make the evidence requirements and their satisfaction
state the interface between declarative programs and physical optimization, and
to evaluate every claim---including the absence of a benefit---against a
matched budget with the same rigor as the presence of one.

\paragraph{Limitations}
The matched-budget headline numbers rest on a \emph{development-split} sample
of matched valid pairs $n=16/17/22$ (HotpotQA / 2WikiMultiHop / MuSiQue,
pooled $n=55$); the $\ge 3$-slot subdomain cost finding is exploratory at
$n=11$ (pooled $1{+}4{+}6$). None of these numbers is validated on a held-out
sealed set, and the pre-registered TKDE protocol for that set is not yet
frozen. All results use a single frozen decoder (qwen3.8-27b at temperature 0)
through a shared provider; generalizability to other decoders is untested.
Evidence corpora are
the public HotpotQA/2Wiki/MuSiQue passage indices; the heterogeneous-evidence
result (HoVer/FEVEROUS pilot) is a small transfer proof-of-concept, not a
Wikipedia-scale retrieval benchmark. The chain law $\tau = 2d - 1$ is exact on
well-defined \emph{chains} where the compiled requirement order respects
topological dependencies; on general (branching) graphs the enumerate-index
proxy is a heuristic and is outside the claim's scope. Cost accounting covers
realized retrieval calls and LLM calls under the matched budget; latency and
monetary cost are reported for the optimizer overhead but are not the primary
comparison axis.
